\documentclass[11pt,a4paper]{article}

% ---------------------------------------------------------------
% Tipografia: Palatino, versalete, uma cor de destaque (bordô)
% ---------------------------------------------------------------
\usepackage[utf8]{inputenc}
\usepackage[T1]{fontenc}
\usepackage[brazil]{babel}
\usepackage[sc,osf]{mathpazo}
\usepackage{microtype}
\usepackage[top=2.4cm, bottom=2.6cm, left=2.6cm, right=2.6cm]{geometry}
\usepackage{amsmath}
\usepackage{amssymb}
\usepackage{xcolor}
\usepackage{tcolorbox}
\tcbuselibrary{skins, breakable}
\usepackage{enumitem}
\usepackage{booktabs}
\usepackage{array}
\usepackage{tabularx}
\usepackage{needspace}
\usepackage{hyperref}

\definecolor{vinho}{RGB}{122, 27, 42}
\definecolor{grafite}{RGB}{70, 70, 70}
\hypersetup{colorlinks=true, linkcolor=vinho, urlcolor=vinho}

\linespread{1.06}
\setlength{\parindent}{0pt}
\setlength{\parskip}{0.45em}

% ---------------------------------------------------------------
% Cabeçalhos de bloco e seção
% ---------------------------------------------------------------
\newcommand{\bloco}[3]{%
  \needspace{7\baselineskip}
  \vspace{1.8em}
  {\scshape\color{vinho}bloco #1}\hfill
    {\small\itshape\color{grafite}\textasciitilde\,#3 min}\\[0.1em]
  {\Large\bfseries #2}\\[-0.55em]
  {\color{black!30}\rule{\linewidth}{0.5pt}}
  \par\nobreak\smallskip
}

\newcommand{\secao}[1]{%
  \needspace{6\baselineskip}
  \vspace{1.8em}
  {\Large\bfseries #1}\\[-0.55em]
  {\color{black!30}\rule{\linewidth}{0.5pt}}
  \par\nobreak\smallskip
}

\newcommand{\objetivo}[1]{{\itshape\color{grafite}Objetivo: #1}\par\smallskip}

% ---------------------------------------------------------------
% Painéis: barra fina à esquerda + rótulo em versalete
% ---------------------------------------------------------------
\newtcolorbox{quadro}{blanker, breakable, left=12pt,
  borderline west={1.4pt}{0pt}{grafite},
  before skip=11pt, after skip=11pt,
  before upper={{\small\scshape\color{grafite}no quadro}\par\smallskip}}

\newtcolorbox{perguntas}{blanker, breakable, left=12pt,
  borderline west={1.4pt}{0pt}{vinho},
  before skip=11pt, after skip=11pt,
  before upper={{\small\scshape\color{vinho}perguntar \`a turma}\par\smallskip}}

\newtcolorbox{obs}{blanker, breakable, left=12pt,
  borderline west={1.4pt}{0pt}{black!30},
  before skip=11pt, after skip=11pt,
  before upper={{\small\scshape\color{black!55}observa\c{c}\~oes}\par\smallskip}}

\newtcolorbox{corte}{blanker, breakable, left=12pt,
  borderline west={1.4pt}{0pt}{vinho!70, dashed},
  before skip=11pt, after skip=11pt,
  before upper={{\small\scshape\color{vinho}se o tempo apertar}\par\smallskip}}

\newtcolorbox{projetor}{blanker, breakable, left=12pt,
  borderline west={1.4pt}{0pt}{grafite},
  borderline west={1.4pt}{3.2pt}{black!25},
  before skip=11pt, after skip=11pt,
  before upper={{\small\scshape\color{grafite}projetor}\par\smallskip}}

% ---------------------------------------------------------------
\begin{document}
% ---------------------------------------------------------------

% ---------------------------------------------------------------
% TÍTULO
% ---------------------------------------------------------------
\thispagestyle{empty}

{\scshape\color{vinho}cursinho solid\'ario \,\textperiodcentered\, f\'isica}\\[-0.5em]
{\color{vinho}\rule{\linewidth}{1.6pt}}\\[0.9em]
{\huge\bfseries Trabalho e Energia}\\[0.35em]
{\LARGE Energia, Conserva\c{c}\~ao, Trabalho e Pot\^encia}\\[0.7em]
{\itshape\color{grafite}Plano de aula \,\textperiodcentered\, 2h
 \,\textperiodcentered\, ENEM e vestibular UTFPR}
\hfill {\color{grafite}2026}\\[-0.3em]
{\color{black!30}\rule{\linewidth}{0.5pt}}

\vspace{1.2em}

% ---------------------------------------------------------------
% VISÃO GERAL
% ---------------------------------------------------------------
{\renewcommand{\arraystretch}{1.45}
\begin{tabularx}{\linewidth}{@{}p{2.6cm}X@{}}
  \multicolumn{2}{@{}l}{{\scshape\color{vinho}vis\~ao geral}}\\
  \addlinespace[2pt]
  \toprule
  \textbf{Objetivo} & Partir da ideia de energia e das suas formas
    mec\^anicas (cin\'etica, potencial gravitacional e el\'astica),
    chegar \`a conserva\c{c}\~ao da energia mec\^anica e s\'o ent\~ao
    definir o trabalho de uma for\c{c}a e a pot\^encia, aplicando tudo
    em quest\~oes de ENEM e do vestibular da UTFPR. \\
  \textbf{Sequ\^encia} & \textbf{Energia primeiro, trabalho depois.} A
    turma j\'a tem intui\c{c}\~ao sobre energia, e trabalho \'e o
    conceito que mais confunde (todo mundo acha que fazer for\c{c}a j\'a
    \'e trabalho). Quando $W = Fd\cos\theta$ aparece, no Bloco 5, ela
    j\'a sabe o que \'e energia e s\'o precisa entender que trabalho \'e
    a energia \textbf{em tr\^ansito}. \\
  \textbf{Formato} & Teoria no quadro do in\'icio ao fim e, no final,
    resolu\c{c}\~ao de quest\~oes tamb\'em no quadro. O projetor s\'o
    é utilizado para mostrar as quest\~oes. \\
  \textbf{Pr\'e-requisito} & Aula de Din\^amica (Leis de Newton, peso,
    normal, atrito e for\c{c}a el\'astica). \\
  \bottomrule
\end{tabularx}}

\vspace{1.4em}

% ---------------------------------------------------------------
% CRONOGRAMA
% ---------------------------------------------------------------
{\renewcommand{\arraystretch}{1.35}
\begin{tabularx}{\linewidth}{@{}clX r@{}}
  \multicolumn{4}{@{}l}{{\scshape\color{vinho}cronograma}}\\
  \addlinespace[2pt]
  \toprule
  \textbf{Bloco} & \textbf{T\'opico} & \textbf{Foco principal} & \textbf{Tempo} \\
  \midrule
  0 & Abertura & Situa\c{c}\~ao-problema: a montanha-russa & 10 min \\
  1 & O que \'e energia & Formas, joule, ``n\~ao some, muda de forma'' & 8 min \\
  2 & Energia cin\'etica & $E_c = mv^2/2$ e o quadrado da velocidade & 10 min \\
  3 & Energia potencial & $E_{pg} = mgh$, $E_{pe} = kx^2/2$, conservativas & 12 min \\
  4 & Conserva\c{c}\~ao & $E_m$ constante e o que muda com atrito & 18 min \\
  5 & Trabalho & $W = Fd\cos\theta$, sinal, peso, teorema, gr\'afico & 20 min \\
  6 & Pot\^encia & $P = W/\Delta t = Fv$, rendimento, kWh & 10 min \\
  7 & Quest\~oes & 8 quest\~oes de vestibular ao vivo (projetor) & 35 min \\
  8 & Encerramento & Resumo das f\'ormulas + tarefa & 7 min \\
  \midrule
  & \textbf{Total} & & \textbf{130 min} \\
  \bottomrule
\end{tabularx}}

% ===================================================================
\bloco{0}{Abertura e Situa\c{c}\~ao-Problema}{10}
% ===================================================================

\objetivo{come\c{c}ar com a montanha-russa, que n\~ao tem motor depois
do primeiro morro, antes de entrar na teoria.}

\begin{perguntas}
Lan\c{c}ar as perguntas oralmente e deixar os alunos responderem, sem
corrigir nada ainda:
\begin{enumerate}[nosep]
  \item ``A montanha-russa \'e puxada por um motor s\'o at\'e o alto do
        primeiro morro. Depois disso o motor desliga.
        \textbf{De onde vem a velocidade no resto do percurso?}''
  \item ``Por que \textbf{nenhum morro depois do primeiro \'e mais
        alto} que ele?''
  \item ``Voc\^e desce a ladeira de bicicleta, freia at\'e parar e o
        freio fica quente. \textbf{Para onde foi a energia do
        movimento?}''
\end{enumerate}
\textit{N\~ao dar as respostas ainda. Anotar as ideias no canto do
quadro. A terceira fecha logo no Bloco 1; a segunda s\'o no Bloco 4.}
\end{perguntas}

\begin{quadro}
\begin{itemize}[nosep]
  \item Escrever o t\'itulo no espa\c{c}o principal:
        \textbf{Energia e Trabalho.}
  \item Desenhar o perfil da montanha-russa, um morro alto e dois
        menores depois, e deixar o desenho no quadro at\'e o Bloco 4.
\end{itemize}
\end{quadro}

\begin{obs}
\textbf{Come\c{c}o da aula:} situar a turma. Na din\^amica a gente
resolvia tudo com $\vec{F}_R = m\vec{a}$ mais cinem\'atica, e precisava
saber a acelera\c{c}\~ao e o tempo. Agora vem um atalho: quando a
quest\~ao liga \textbf{velocidade} com \textbf{altura} ou
\textbf{dist\^ancia}, e n\~ao fala em tempo, energia resolve bem mais
r\'apido. \'E o assunto de mec\^anica que mais cai no ENEM.

\textbf{Demonstra\c{c}\~ao:} mostrar vídeo do professor Walter Lewin (https://www.youtube.com/watch?v=77ZF50ve6rs)
e mostrar que assim como na montanha-russa, o pêndulo nunca atinge uma altura maior que a inicial,
se não houver forças agindo nele.
\end{obs}

% ===================================================================
\bloco{1}{O Que \'e Energia}{8}
% ===================================================================

\objetivo{dar sentido \`a palavra ``energia'' antes de qualquer
f\'ormula e instalar a regra que sustenta a aula inteira: a energia
muda de forma e o total n\~ao muda.}

\begin{quadro}
\textbf{1. Uma defini\c{c}\~ao honesta} (2 min)
\begin{itemize}[nosep]
  \item {Energia \'e o que um sistema
        precisa ter para provocar mudan\c{c}as}, p\^or algo em
        movimento, levantar, deformar, aquecer, iluminar.
  \item Ser honesto com a turma: ningu\'em define energia em uma frase
        s\'o. O que a f\'isica faz \'e ensinar a \textbf{calcular} um
        n\'umero em cada situa\c{c}\~ao. A gra\c{c}a \'e que a soma
        desses n\'umeros n\~ao muda.
  \item Unidade: \textbf{joule} (J), a mesma para todas as formas. \'E
        justamente isso que permite somar coisas diferentes e
        compar\'a-las.
\end{itemize}

\medskip
\textbf{2. As formas de energia} (3 min)
\begin{itemize}[nosep]
  \item Listar no quadro: cin\'etica (movimento), potencial
        gravitacional (altura), potencial el\'astica
        (deforma\c{c}\~ao), t\'ermica (calor), el\'etrica, qu\'imica
        (comida, combust\'ivel, pilha), nuclear.
  \item \textbf{Circular as tr\^es primeiras}: s\~ao as da aula de hoje
        e, juntas, formam a \textbf{energia mec\^anica}.
  \item Montar uma cadeia com setas, para mostrar que \'e sempre a
        mesma energia trocando de roupa:
        \[
          \underbrace{\text{comida}}_{\text{qu\'imica}} \rightarrow
          \underbrace{\text{pedalar}}_{\text{cin\'etica}} \rightarrow
          \underbrace{\text{subir o morro}}_{\text{potencial}} \rightarrow
          \underbrace{\text{frear}}_{\text{t\'ermica}}
        \]
\end{itemize}

\medskip
\textbf{3. A regra do jogo} (3 min)
\begin{itemize}[nosep]
  \item \textbf{A energia n\~ao some, ela muda de forma}.
  \item \textbf{Responder a terceira pergunta da abertura:} a energia
        do movimento virou calor no freio, no pneu e no ar. Ela n\~ao
        sumiu, ficou \textbf{espalhada}, e espalhada n\~ao d\'a mais
        para aproveitar.
  \item Introduzir a palavra \textbf{degradar}: a gente n\~ao ``gasta''
        energia, a gente degrada, transforma energia \'util em calor
        espalhado.
\end{itemize}
\end{quadro}

\begin{perguntas}
\begin{itemize}[nosep]
  \item ``Uma pilha nova e uma pilha descarregada pesam praticamente o
        mesmo. O que acabou na descarregada?''
        \textit{(A energia qu\'imica guardada nas subst\^ancias de
        dentro: elas viraram outras, com menos energia guardada.)}
  \item ``De onde vem, no fundo, a energia da gasolina?''
        \textit{(Do Sol. Plantas guardaram energia solar como energia
        qu\'imica, h\'a milh\~oes de anos.)}
\end{itemize}
\end{perguntas}

\begin{obs}
Este bloco \'e curto de prop\'osito. O objetivo \'e s\'o instalar a
palavra e a regra; a precis\~ao vem nos blocos seguintes, quando
aparecem as f\'ormulas. \textbf{N\~ao deixar virar aula de ``tipos de
energia''}: 8 minutos e segue.
\end{obs}

\begin{corte}
Cortar a cadeia de transforma\c{c}\~oes (item 2) e as perguntas. N\~ao
cortar a frase ``a energia n\~ao some, ela muda de forma'': a aula
inteira se apoia nela.
\end{corte}

% ===================================================================
\bloco{2}{Energia Cin\'etica}{10}
% ===================================================================

\objetivo{dar a primeira f\'ormula da aula e martelar que a velocidade
entra ao quadrado, que \'e o que o ENEM mais cobra nesse tema.}

\begin{quadro}
\textbf{1. A energia do movimento} (4 min)
\begin{itemize}[nosep]
  \item Escrever na faixa de resumo:
  \[
    \boxed{E_c = \frac{m\,v^2}{2}}
  \]
  medida em joules e \textbf{nunca negativa}, porque a velocidade
  est\'a ao quadrado.
  \item Todo corpo em movimento tem energia cin\'etica. Parado,
        $E_c = 0$.
\end{itemize}

\medskip
\textbf{2. O quadrado da velocidade} (4 min)
\begin{itemize}[nosep]
  \item \textbf{Insistir}: se a velocidade dobra, a energia cin\'etica
        fica quatro vezes maior; se triplica, fica nove vezes maior.
  \item Montar a tabela no quadro, com um carro de 1\,000 kg:
  \begin{center}\small
    \begin{tabular}{@{}lll@{}}
      \toprule
      $v$ & $E_c$ & compara\c{c}\~ao \\
      \midrule
      10 m/s (36 km/h) & 50 kJ  & --- \\
      20 m/s (72 km/h) & 200 kJ & 4 vezes \\
      30 m/s (108 km/h) & 450 kJ & 9 vezes \\
      \bottomrule
    \end{tabular}
  \end{center}
  \item Consequ\^encia pr\'atica: a dist\^ancia de frenagem de um carro
        cresce bem mais r\'apido do que a velocidade. \'E o gancho que
        o ENEM usa nas quest\~oes de seguran\c{c}a no tr\^ansito.
\end{itemize}

\medskip
\textbf{3. Onde ela aparece} (2 min)
\begin{itemize}[nosep]
  \item Voltar ao desenho da montanha-russa: no ponto mais baixo \'e
        onde a energia cin\'etica \'e m\'axima.
  \item Ainda n\~ao d\'a para calcular quanto vale --- falta a outra
        metade da hist\'oria, a energia potencial, que vem agora.
\end{itemize}
\end{quadro}

\begin{perguntas}
\begin{itemize}[nosep]
  \item ``Um caminh\~ao e uma moto andam na mesma velocidade. Qual tem
        mais energia cin\'etica?''
        \textit{(O caminh\~ao, a massa entra direto na f\'ormula.)}
  \item ``E se a moto estiver ao dobro da velocidade do caminh\~ao?''
        \textit{(A\'i depende dos valores. Serve para mostrar que
        velocidade entra ao quadrado e massa n\~ao.)}
\end{itemize}
\end{perguntas}

\begin{obs}
Se algu\'em perguntar ``de onde sai esse $mv^2/2$?'', responder que ele
vem da din\^amica e que a dedu\c{c}\~ao aparece no \textbf{Bloco 5},
junto com o teorema. N\~ao deduzir agora: sem trabalho definido, a
conta n\~ao fecha.
\end{obs}

\begin{corte}
Cortar a tabela num\'erica do item 2 e o item 3 inteiro. N\~ao cortar a
conversa sobre o quadrado da velocidade.
\end{corte}

% ===================================================================
\bloco{3}{Energia Potencial}{12}
% ===================================================================

\objetivo{apresentar as duas energias guardadas e separar for\c{c}a
conservativa de dissipativa, que \'e o crit\'erio para saber se a
energia mec\^anica se conserva.}

\begin{quadro}
\textbf{1. Energia potencial gravitacional} (5 min)
\begin{itemize}[nosep]
  \item Escrever na faixa de resumo:
  \[
    \boxed{E_{pg} = m\,g\,h}
  \]
  \item \'E energia de \textbf{posi\c{c}\~ao}. O corpo l\'a em cima
        ainda n\~ao tem velocidade, mas tem a promessa dela.
  \item \textbf{O n\'ivel zero \'e escolha nossa.} O valor de $E_{pg}$
        muda, mas a \textbf{varia\c{c}\~ao} n\~ao. Combinar com a turma
        de escolher sempre o ponto mais baixo do problema como zero.
  \item Voltar \`a montanha-russa: no alto do primeiro morro \'e tudo
        $E_{pg}$, e a\'i est\'a a resposta que ficou pendente no Bloco
        2 sobre de onde vem a velocidade l\'a embaixo.
\end{itemize}

\medskip
\textbf{2. Energia potencial el\'astica} (4 min)
\begin{itemize}[nosep]
  \item Retomar a Lei de Hooke: $F = k\,x$, com
        $x$ a \textbf{deforma\c{c}\~ao}, e n\~ao o comprimento da mola.
  \item
  \[
    \boxed{E_{pe} = \frac{k\,x^2}{2}}
  \]
  \item Vale para mola comprimida e para mola esticada, porque $x$
        est\'a ao quadrado. Exemplos: estilingue, arco e flecha,
        amortecedor, trampolim.
  \item Mesma pegadinha da cin\'etica: comprimir at\'e a metade guarda
        \textbf{um quarto} da energia, e n\~ao a metade.
  \item Deixar uma promessa no quadro: essa f\'ormula sai da
        \textbf{\'area de um gr\'afico}. A gente confirma isso no Bloco
        5, quando trabalho estiver definido.
\end{itemize}

\medskip
\textbf{3. Conservativas e dissipativas} (3 min)
\begin{center}\small
  \begin{tabular}{@{}p{4.2cm}p{7.2cm}@{}}
    \toprule
    \textbf{Conservativas} & peso e for\c{c}a el\'astica. Guardam a
      energia e \textbf{devolvem} depois; t\^em energia potencial
      associada \\
    \textbf{Dissipativas} & atrito e resist\^encia do ar. N\~ao
      devolvem, a energia vira \textbf{calor} \\
    \bottomrule
  \end{tabular}
\end{center}
Escrever no quadro o crit\'erio que abre o pr\'oximo bloco: \textbf{se
s\'o houver for\c{c}as conservativas, a energia mec\^anica se
conserva}.
\end{quadro}

\begin{perguntas}
\begin{itemize}[nosep]
  \item ``Uma pessoa no quinto andar tem energia potencial. Medida a
        partir de onde?''
        \textit{(Do que a gente escolher. O que tem sentido f\'isico \'e
        a diferen\c{c}a entre dois pontos.)}
  \item ``Se eu comprimir a mola s\'o at\'e a metade, guardo metade da
        energia?''
        \textit{(N\~ao, um quarto. \'E a mesma pegadinha do quadrado da
        energia cin\'etica.)}
\end{itemize}
\end{perguntas}

\begin{corte}
Cortar a energia el\'astica (item 2) e avisar que ela volta junto com a
Lei de Hooke. N\~ao cortar a conversa sobre o n\'ivel de refer\^encia
nem a tabela de conservativa e dissipativa, sem elas o Bloco 4 n\~ao
fecha.
\end{corte}

% ===================================================================
\bloco{4}{Conserva\c{c}\~ao da Energia Mec\^anica}{18}
% ===================================================================

\objetivo{chegar no resultado central da aula, fechar a
situa\c{c}\~ao-problema da abertura e mostrar o que muda quando existe
atrito. \'E o bloco mais importante para o ENEM.}

\begin{quadro}
\textbf{1. Energia mec\^anica} (4 min)
\begin{itemize}[nosep]
  \item Definir: $E_m = E_c + E_p$, a cin\'etica mais todas as
        potenciais. \'E a soma das tr\^es formas que foram circuladas
        no Bloco 1.
  \item Escrever:
  \[
    \boxed{\text{sem atrito:}\quad
      E_{c_i} + E_{p_i} = E_{c_f} + E_{p_f}}
  \]
  \item A energia \textbf{troca de forma},
        potencial vira cin\'etica e cin\'etica vira potencial, mas o
        \textbf{total n\~ao muda}. \'E a regra do Bloco 1 aplicada s\'o
        \`a parte mec\^anica.
\end{itemize}

\medskip
\textbf{2. Voltar na montanha-russa} (7 min)
\begin{itemize}[nosep]
  \item Usar o desenho da montanha-russa. Marcar
        tr\^es pontos: o alto do primeiro morro, o vale e o alto do
        segundo morro.
  \item Em cada ponto escrever quanto vale a cin\'etica e quanto vale a
        potencial, com barrinhas proporcionais, que fica visual. No
        topo \'e tudo potencial, no vale \'e tudo cin\'etica, no
        segundo morro fica dividido.
  \item \textbf{Responder as duas perguntas da abertura:} a velocidade
        vem da energia potencial do primeiro morro, e nenhum morro pode
        ser mais alto que ele porque \textbf{n\~ao sobra energia} para
        chegar l\'a em cima.
  \item Comentar que num escorregador liso vale a mesma conta. A
        trajet\'oria n\~ao importa, s\'o a altura.
\end{itemize}

\medskip
\textbf{3. Quando tem atrito} (7 min)
\begin{itemize}[nosep]
  \item A energia mec\^anica \textbf{n\~ao} se conserva: parte dela
        vira calor. Escrever na faixa de resumo:
  \[
    \boxed{E_{m_i} = E_{m_f} + E_{\text{dissipada}}}
  \]
  \item Deixar claro que \textbf{a energia total continua se
        conservando}, s\'o a mec\^anica que diminui. Nada desaparece
        --- \'e a regra do Bloco 1 outra vez.
  \item Por enquanto, $E_{\text{dissipada}}$ \'e s\'o ``o que faltou''
        na conta. A f\'ormula que d\'a o valor dela,
        $f_{at}\cdot d$, sai no \textbf{Bloco 5}, quando trabalho
        estiver definido. Avisar isso, para ningu\'em achar que ficou
        buraco.
  \item Exemplos do dia a dia: o freio do carro esquenta, a m\~ao
        esquenta quando a gente esfrega uma na outra, o meteoro pega
        fogo ao entrar na atmosfera.
  \item Citar o \textbf{freio regenerativo} do carro el\'etrico. Em vez
        de jogar a energia fora como calor, o motor vira gerador e
        devolve parte dela para a bateria. \'E contexto de ENEM.
\end{itemize}
\end{quadro}

\begin{perguntas}
\begin{itemize}[nosep]
  \item ``Dois escorregadores da mesma altura, um reto e um cheio de
        curvas, os dois sem atrito. Em qual a crian\c{c}a chega embaixo
        com maior velocidade? E em qual ela chega mais r\'apido?''
        \textit{(A velocidade final \'e a mesma nos dois, s\'o depende
        da altura. O tempo \'e que pode ser diferente. Separar bem as
        duas perguntas, \'e onde a turma se perde.)}
  \item ``Se a energia n\~ao se cria nem se destr\'oi, por que a gente
        fala em economizar energia?''
        \textit{(Retomar a palavra do Bloco 1: a gente n\~ao gasta
        energia, a gente \textbf{degrada}. Transforma energia \'util em
        calor espalhado, que n\~ao d\'a mais para aproveitar.)}
\end{itemize}
\end{perguntas}

\begin{obs}
Armadilha cl\'assica do ENEM: usar conserva\c{c}\~ao de energia num
problema que tem atrito. Antes de escrever $E_{m_i} = E_{m_f}$, obrigar
a turma a fazer a pergunta: \textbf{tem atrito ou resist\^encia do ar
nessa quest\~ao?} Se o enunciado falar em ``superf\'icie lisa'',
``atrito desprez\'ivel'' ou ``sem resist\^encia do ar'', pode
conservar.
\end{obs}

\begin{corte}
Cortar o freio regenerativo e as barrinhas do item 2, deixando s\'o a
fala em cima do desenho. N\~ao cortar o item 3.
\end{corte}

% ===================================================================
\bloco{5}{Trabalho de uma For\c{c}a}{20}
% ===================================================================

\objetivo{apresentar o trabalho como a medida de \textbf{quanta energia
uma for\c{c}a transfere} e, com ele, fechar as duas contas deixadas em
aberto de prop\'osito: o teorema $W = \Delta E_c$ e o valor da energia
dissipada pelo atrito.}

\begin{perguntas}
Abrir o bloco com a pergunta, antes de qualquer f\'ormula, e deixar no
ar por uns 30 segundos:
\begin{itemize}[nosep]
  \item ``Empurrar uma parede com toda a for\c{c}a durante 10 minutos:
        voc\^e realizou \textbf{trabalho}?''
\end{itemize}
\end{perguntas}

\begin{quadro}
\textbf{1. Defini\c{c}\~ao} (5 min)
\begin{itemize}[nosep]
  \item Escrever na faixa de resumo:
  \[
    \boxed{W = F \cdot d \cdot \cos\theta}
  \]
  em que $\theta$ \'e o \^angulo entre a \textbf{for\c{c}a} e o
  \textbf{deslocamento}.
  \item Unidade: joule, $1\ \text{J} = 1\ \text{N}\cdot\text{m}$.
        \textbf{A mesma unidade da energia}, e isso n\~ao \'e
        coincid\^encia.
  \item Escrever a frase no quadro: \textbf{trabalho \'e energia
        transferida por uma for\c{c}a}. Trabalho positivo entrega
        energia ao corpo, negativo tira energia dele. Trabalho n\~ao
        \'e um tipo de energia, \'e energia \textbf{em tr\^ansito}.
  \item Responder a pergunta da parede: a parede n\~ao sai do lugar,
        ent\~ao $d = 0$ e $W = 0$. Nenhuma energia foi transferida para
        ela. Cansa\c{c}o n\~ao \'e trabalho em f\'isica.
\end{itemize}

\medskip
\textbf{2. Os tr\^es sinais e o atrito} (5 min)
\begin{itemize}[nosep]
  \item Montar a tabela no quadro:
  \begin{center}\small
    \begin{tabular}{@{}lll@{}}
      \toprule
      \^Angulo & Sinal & Nome \\
      \midrule
      $\theta < 90^\circ$ & $W > 0$ & motor (ajuda o movimento) \\
      $\theta = 90^\circ$ & $W = 0$ & nulo \\
      $\theta > 90^\circ$ & $W < 0$ & resistente (freia) \\
      \bottomrule
    \end{tabular}
  \end{center}
  \item \textbf{Atrito cin\'etico \'e sempre resistente}, com
        $\theta = 180^\circ$, ent\~ao $W_{at} = -f_{at}\,d$.
  \item \textbf{Fechar o buraco do Bloco 4:} em m\'odulo, esse \'e
        exatamente o valor que faltava.
  \[
    \boxed{E_{\text{dissipada}} = |W_{at}| = f_{at}\cdot d}
  \]
  \item \textbf{For\c{c}a perpendicular ao deslocamento n\~ao realiza
        trabalho.} Dois casos que caem direto: a normal num corpo que
        desliza na horizontal e a for\c{c}a centr\'ipeta no movimento
        circular.
  \item Exemplo: carregar uma mochila andando na horizontal. A
        for\c{c}a que segura a mochila \'e vertical e o deslocamento \'e
        horizontal, ent\~ao o trabalho \'e zero.
\end{itemize}

\medskip
\textbf{3. Trabalho do peso} (4 min)
\begin{itemize}[nosep]
  \item Escrever na faixa de resumo:
  \[
    W_P = \pm\, m\,g\,h
  \]
  positivo \textbf{descendo}, negativo \textbf{subindo}, com $h$ o
  desn\'ivel \textbf{vertical}.
  \item O ponto principal: \textbf{o trabalho do peso n\~ao depende da
        trajet\'oria}, s\'o do desn\'ivel. Desenhar tr\^es caminhos
        entre os mesmos dois pontos (queda vertical, rampa e uma curva
        qualquer) e escrever o mesmo $mgh$ nos tr\^es.
  \item Num percurso fechado, que sai e volta ao mesmo ponto, o
        trabalho do peso \'e \textbf{zero}. \'E exatamente por isso que
        o peso foi chamado de \textbf{conservativo} no Bloco 3.
  \item Ligar com o Bloco 3: $W_P = -\Delta E_{pg}$. O peso faz
        trabalho positivo justamente quando a energia potencial
        diminui. \'E a mesma coisa dita de duas maneiras.
\end{itemize}

\medskip
\textbf{4. Teorema do trabalho e energia cin\'etica} (4 min)
\begin{itemize}[nosep]
  \item Escrever na faixa de resumo:
  \[
    \boxed{W_{\text{total}} = \Delta E_c = \frac{m v_f^2}{2}
      - \frac{m v_i^2}{2}}
  \]
  \item Ler em portugu\^es: \textbf{o trabalho da resultante \'e o que
        muda a velocidade do corpo}. Se o trabalho total \'e zero, a
        velocidade n\~ao muda --- \'e a 1\textordfeminine{} Lei escrita
        em energia.
  \item Destacar o ``total'': \'e a soma do trabalho de
        \textbf{todas} as for\c{c}as, ou o trabalho da resultante, que
        d\'a no mesmo.
  \item \textbf{De onde vem o $mv^2/2$} (a pergunta que ficou no Bloco
        2): pegar $v_f^2 = v_i^2 + 2a\Delta S$, isolar $a\Delta S$ e
        multiplicar tudo por $m$. Aparece $F_R\,\Delta S$ de um lado e
        a diferen\c{c}a das energias cin\'eticas do outro.
  \item Avisar que os problemas de frenagem e de arrancada que a gente
        fazia com Torricelli mais $F = ma$ saem em uma conta s\'o por
        aqui. Isso fica claro no Bloco 7.
\end{itemize}

\medskip
\textbf{5. \'Area do gr\'afico $F \times d$} (2 min)
\begin{itemize}[nosep]
  \item Quando a for\c{c}a \textbf{varia}, n\~ao d\'a para usar
        $W = Fd\cos\theta$ direto. Vale sempre:
        \textbf{o trabalho \'e a \'area embaixo do gr\'afico
        $F \times d$}.
  \item Desenhar dois gr\'aficos lado a lado: for\c{c}a constante, que
        d\'a um ret\^angulo, e for\c{c}a de mola crescendo com $x$, que
        d\'a um tri\^angulo de \'area $kx^2/2$. \textbf{\'E a promessa
        do Bloco 3 cumprida}: a energia el\'astica n\~ao \'e f\'ormula
        nova, \'e essa \'area.
  \item \'Area abaixo do eixo conta como trabalho negativo.
\end{itemize}
\end{quadro}

\begin{perguntas}
\begin{itemize}[nosep]
  \item ``Subir at\'e o segundo andar pela escada ou por uma rampa bem
        longa: em qual dos dois o \textbf{trabalho contra o peso} \'e
        maior?''
        \textit{(Igual nos dois, $mgh$. A rampa pede menos
        \textbf{for\c{c}a}, mas por uma dist\^ancia maior. \'E para
        isso que a rampa serve.)}
  \item ``A Lua d\'a voltas em torno da Terra h\'a bilh\~oes de anos. A
        gravidade da Terra realiza trabalho sobre ela?''
        \textit{(Numa \'orbita circular n\~ao, a for\c{c}a \'e
        perpendicular ao deslocamento. Por isso a velocidade n\~ao
        muda.)}
\end{itemize}
\end{perguntas}

\begin{obs}
Confus\~ao mais comum: achar que fazer for\c{c}a j\'a \'e trabalho. Sem
deslocamento \textbf{na dire\c{c}\~ao da for\c{c}a} o trabalho \'e
zero. Repetir isso nos tr\^es exemplos, parede, mochila e for\c{c}a
centr\'ipeta, at\'e colar.

Vantagem de dar trabalho depois de energia: a turma j\'a sabe o que \'e
energia, ent\~ao d\'a para apresentar trabalho como \textbf{a conta de
quanta energia mudou de m\~aos}, e n\~ao como mais uma f\'ormula solta.
\end{obs}

\begin{corte}
Cortar o gr\'afico (item 5) e a pergunta da Lua. \textbf{N\~ao cortar}
o teorema (item 4) nem a f\'ormula da energia dissipada (item 2): os
dois fecham buracos deixados de prop\'osito nos blocos anteriores.
\end{corte}

% ===================================================================
\bloco{6}{Pot\^encia e Rendimento}{10}
% ===================================================================

\objetivo{fechar a teoria com pot\^encia, que \'e o que aparece nas
quest\~oes de usina, motor e conta de luz, e separar energia de
pot\^encia.}

\begin{quadro}
\textbf{1. Pot\^encia} (5 min)
\begin{itemize}[nosep]
  \item
  \[
    \boxed{P = \frac{W}{\Delta t} = \frac{\Delta E}{\Delta t}}
    \qquad
    1\ \text{W} = 1\ \text{J/s}
  \]
  \item Pot\^encia diz quanta energia passa por
        segundo. A mesma energia em menos tempo d\'a mais pot\^encia.
  \item Comparar duas pessoas subindo a mesma escada, uma correndo e
        outra devagar. As duas gastam a mesma energia, mas a que sobe
        correndo tem mais pot\^encia.
  \item Caso particular, quando a velocidade \'e constante:
        $P = F\cdot v$. \'E o que resolve quest\~ao de motor de carro e
        de esteira.
\end{itemize}

\medskip
\textbf{2. Rendimento} (3 min)
\begin{itemize}[nosep]
  \item $\eta = \dfrac{P_{\text{\'util}}}{P_{\text{total}}}$, sempre
        menor que 100\%.
  \item O resto vira calor, som e vibra\c{c}\~ao. \textbf{Nenhuma
        m\'aquina \'e 100\% eficiente}, e alternativa que fala em
        m\'aquina de movimento perp\'etuo est\'a errada. \'E a
        degrada\c{c}\~ao do Bloco 1 aparecendo em n\'umero.
  \item Formato cl\'assico de quest\~ao: a usina hidrel\'etrica. A
        \'agua cai de uma altura $h$, o $mgh$ por segundo d\'a a
        pot\^encia bruta e o rendimento da turbina d\'a a pot\^encia
        el\'etrica.
\end{itemize}

\medskip
\textbf{3. Quilowatt-hora} (2 min)
\begin{itemize}[nosep]
  \item A conta de luz cobra \textbf{energia}, n\~ao pot\^encia. Um kWh
        \'e a energia de um aparelho de 1000 W ligado por uma hora, ou
        seja, $3{,}6\times 10^6$ J.
  \item Chuveiro e ar-condicionado pesam na conta porque t\^em
        pot\^encia alta, mesmo ficando pouco tempo ligados.
\end{itemize}
\end{quadro}

\begin{perguntas}
\begin{itemize}[nosep]
  \item ``Um elevador sobe dez andares devagar e outro sobe os mesmos
        dez andares r\'apido. Qual gastou mais \textbf{energia}? Qual
        tem mais \textbf{pot\^encia}?''
        \textit{(A energia \'e a mesma, o $mgh$ \'e o mesmo. A
        pot\^encia \'e maior no que subiu r\'apido. Essa
        distin\c{c}\~ao \'e o objetivo do bloco.)}
\end{itemize}
\end{perguntas}

\begin{corte}
Cortar o rendimento e o quilowatt-hora (itens 2 e 3), deixando s\'o
$P = W/\Delta t$. Se cortar, avisar que os dois caem em quest\~ao de
usina e passar como tarefa.
\end{corte}

% ===================================================================
\bloco{7}{Resolu\c{c}\~ao de Quest\~oes}{35}
% ===================================================================

\objetivo{terminar a aula resolvendo quest\~oes de verdade,
projetadas uma por uma e resolvidas ao vivo no quadro. O roteiro
\'e puxado para \textbf{ENEM e UTFPR}, que s\~ao as provas que a
turma vai prestar: seis das dez do roteiro e as duas de b\^onus
s\~ao dessas duas bancas. Projetar o arquivo
\texttt{slides\_questoes\_\allowbreak trabalho\_\allowbreak energia.pdf},
que tem uma quest\~ao por slide mais o slide de gabarito, na mesma
ordem desta tabela. Na lista impressa elas s\~ao as marcadas com
{\scshape\color{vinho}[fazer em sala]}.}

\begin{projetor}
Projetar o arquivo
\texttt{slides\_questoes\_\allowbreak trabalho\_\allowbreak energia.pdf}
(um slide por quest\~ao, seguido do slide de gabarito). A lista
impressa,
\texttt{lista\_questoes\_\allowbreak trabalho\_\allowbreak energia.pdf},
fica com a turma.

\smallskip
\end{projetor}

\begin{quadro}
\textbf{Roteiro de quest\~oes:}

\medskip
{\small
\begin{tabular}{@{}cllr@{}}
  \toprule
  \# & Cobra & Fonte & Tempo \\
  \midrule
  Q8  & Conserva\c{c}\~ao: salto com vara           & ENEM 2011          & 3 min \\
  Q11 & Energia el\'astica: mola do brinquedo       & ENEM 2018          & 4 min \\
  Q12 & Barreira de pneus: for\c{c}a e energia      & ENEM 2022          & 3 min \\
  Q15 & Pot\^encia do guindaste ($P = Fv$)          & UTFPR Ver.\ 2026   & 2 min \\
  Q17 & Rendimento e kWh (fotovoltaico)             & UTFPR Inv.\ 2023   & 5 min \\
  Q16 & Pot\^encia n\~ao aproveitada em Itaipu      & ENEM 2016          & 4 min \\[0.4em]
  Q10 & Conserva\c{c}\~ao: bola no ar (num\'erica)  & UERJ               & 4 min \\
  Q13 & Energia dissipada na queda (num\'erica)     & UNIFOR CE          & 4 min \\
  Q1  & Trabalho e o \^angulo da for\c{c}a          & UFF                & 2 min \\
  Q5  & Teorema $W = \Delta E_c$ (num\'erica)       & FGV                & 4 min \\
  \bottomrule
\end{tabular}}

\medskip
\textbf{B\^onus}, se sobrar tempo: \textbf{Q18} (UTFPR Ver.\ 2025,
2 min --- a mesma conta de kWh da Q17, s\'o que sem rendimento; serve
para mostrar que a UTFPR repete o tema fotovoltaico prova ap\'os
prova) e \textbf{Q19} (ENEM 2015, carro solar, 5 min --- fecha a
sequ\^encia: painel, rendimento, pot\^encia e energia cin\'etica na
mesma quest\~ao).

\medskip
A ordem \'e a da aula, e \textbf{n\~ao} a ordem num\'erica da lista.
As \textbf{seis primeiras s\~ao ENEM e UTFPR}, na sequ\^encia dos
temas: conserva\c{c}\~ao (Q8), energia el\'astica (Q11),
dissipa\c{c}\~ao (Q12) e depois o bloco de pot\^encia e rendimento
(Q15, Q17 e Q16), que \'e o formato preferido das duas bancas nesse
tema e quase sempre o que a turma erra por causa da unidade. Se a
aula atrasar, o corte cai nas quatro \'ultimas, e o que a turma vai
ver na prova j\'a foi.

\medskip
As \textbf{quatro \'ultimas} s\~ao de outras bancas e entram como
treino com conta: Q10 e Q13 repetem conserva\c{c}\~ao e
dissipa\c{c}\~ao com n\'umero (o ENEM cobra os dois no conceito, e
essas fecham a m\~ao), Q1 corrige o erro de sinal do trabalho e Q5 \'e
a ponte com a din\^amica.

\medskip
Na lista, essas quest\~oes est\~ao marcadas com
{\scshape\color{vinho}[fazer em sala]} e as duas de b\^onus com
{\scshape\color{vinho}[sala --- b\^onus]}. O mapa da segunda p\'agina
da lista mostra tudo de uma vez. As de outras bancas (Q2, Q3, Q4, Q6,
Q7, Q9 e Q14) ficam de tarefa: treinam a mesma ideia, mas s\~ao
enunciados mais secos que os do ENEM.
\end{quadro}

\begin{obs}
Resolver sempre com a mesma estrutura de tr\^es linhas no quadro:
estado inicial, estado final e, se tiver atrito, descontar o que foi
dissipado. A turma precisa sair da aula com um \textbf{procedimento},
e n\~ao com seis f\'ormulas soltas.

Em alguma quest\~ao num\'erica, se der tempo, resolver de novo pela
din\^amica, com $F = ma$ e Torricelli, e comparar o tamanho das duas
solu\c{c}\~oes. \'E o que convence a turma a usar energia.
\end{obs}

\begin{corte}
Menos de 25 min? Fazer Q8, Q11, Q13, Q15, Q17 e Q5, que cobrem
conserva\c{c}\~ao, energia el\'astica, dissipa\c{c}\~ao, pot\^encia,
rendimento e o teorema. Menos de 15 min? Q8, Q13 e Q17, que s\~ao os
tr\^es t\'opicos que mais caem no ENEM e na UTFPR. Menos de 10?
Q8 e Q17.
\end{corte}

% ===================================================================
\bloco{8}{Encerramento e Tarefa}{7}
% ===================================================================

\begin{quadro}
\textbf{1. Resumo} (4 min). Fechar a faixa lateral que ficou a aula
toda no quadro, lendo tudo de uma vez e \textbf{na ordem em que
apareceu}:
\begin{center}\small
  \begin{tabular}{@{}ll@{}}
    \toprule
    $E_c = mv^2/2$ & energia do movimento \\
    $E_{pg} = mgh$ \quad $E_{pe} = kx^2/2$ & energia guardada \\
    $E_{m_i} = E_{m_f} + E_{\text{dissipada}}$ & conserva\c{c}\~ao \\
    $W = Fd\cos\theta$ & trabalho, energia transferida \\
    $W_{\text{total}} = \Delta E_c$ & teorema do trabalho e energia \\
    $P = W/\Delta t = Fv$ & pot\^encia \\
    \bottomrule
  \end{tabular}
\end{center}

\medskip
\textbf{2. Quando usar energia no lugar de $F = ma$} (3 min).
Escrever o crit\'erio no quadro, \'e o que mais ajuda na prova:
\begin{itemize}[nosep]
  \item A quest\~ao liga \textbf{velocidade com altura ou
        dist\^ancia} e \textbf{n\~ao fala em tempo}: usar
        \textbf{energia}.
  \item A quest\~ao pede \textbf{for\c{c}a} ou
        \textbf{acelera\c{c}\~ao}, ou envolve \textbf{tempo}: usar
        \textbf{din\^amica}.
  \item Trajet\'oria curva ou complicada, sem atrito: usar
        \textbf{energia}, porque a trajet\'oria n\~ao importa.
\end{itemize}
\end{quadro}

\begin{obs}
Passar como tarefa as quest\~oes que n\~ao foram feitas em sala e
avisar qual \'e a pr\'oxima aula.
\end{obs}

% ===================================================================
\secao{Adapta\c{c}\~ao por Tempo Dispon\'ivel}
% ===================================================================

{\renewcommand{\arraystretch}{1.5}
\begin{tabularx}{\linewidth}{@{}lX@{}}
  \toprule
  \textbf{Tempo} & \textbf{O que fazer} \\
  \midrule
  $\approx 1$h30 &
    Bloco0 (5 min) + Bloco1 (5 min) + Bloco2 sem a tabela (7 min) +
    Bloco3 sem a el\'astica (8 min) + Bloco4 completo (18 min) +
    Bloco5 sem o gr\'afico (16 min) + Bloco6 s\'o $P = W/\Delta t$
    (5 min) + 6 quest\~oes (22 min) + fechamento (4 min) \\
  $\approx 2$h &
    Plano completo. Se atrasar, cortar do \textbf{fim} do roteiro
    (Q5 e Q1), porque o bloco ENEM/UTFPR vem primeiro de prop\'osito,
    e segurar o ritmo nas conceituais \\
  $\approx 2$h30 &
    Plano completo + resolver uma quest\~ao num\'erica pelos dois
    m\'etodos, energia e $F = ma$, + as duas de b\^onus (Q18 e Q19) +
    a Q14, que junta mola, atrito e rampa \\
  \bottomrule
\end{tabularx}}

\bigskip

{\scshape\color{vinho}prioridade dos t\'opicos para enem/utfpr}
\begin{enumerate}[nosep]
  \item \textbf{Conserva\c{c}\~ao da energia mec\^anica} (montanha-russa,
        rampa, p\^endulo, queda), que \'e o que mais cai no ENEM
  \item \textbf{Energia dissipada pelo atrito} e a ideia de que a
        energia se degrada, e n\~ao some
  \item \textbf{Pot\^encia e rendimento} (usina, motor, chuveiro, kWh),
        muito frequente em quest\~ao de contexto
  \item \textbf{Teorema $W = \Delta E_c$} e o quadrado da velocidade
        (frenagem, seguran\c{c}a no tr\^ansito)
  \item \textbf{Sinal do trabalho} e for\c{c}a perpendicular, que pega
        muita gente
  \item \textbf{Energia el\'astica} ($kx^2/2$) junto com a Lei de
        Hooke, estilo UTFPR
\end{enumerate}
\end{document}
